\chapter{Installation de ROS2}

\noindent Il est possible de retrouver ce tutoriel en cliquant \href{https://docs.ros.org/en/humble/Installation.html}{ici}. Il est conseillé de suivre le tutoriel sur le site internet plutôt que sur le présent document car il est plus simple de \textit{copier-coller} les commandes (les retours à la ligne, symbolisés par des flèches sur ce document, posent problème).

\section{Configuration du système}

Avant d'installer ROS2, il est nécessaire de mettre à jour le système. Suivre les étapes suivantes :

\begin{enumerate}
\item Mettre à jour les paquets et installer les locales :
\begin{lstlisting}[style=commandstyle]
sudo apt update && sudo apt install locales
\end{lstlisting}

\item Générer les locales pour l'anglais des États-Unis (\texttt{en\_US}) :
\begin{lstlisting}[style=commandstyle]
sudo locale-gen en_US en_US.UTF-8
sudo update-locale LC_ALL=en_US.UTF-8 LANG=en_US.UTF-8
\end{lstlisting}

\item Exporter la variable d'environnement pour la langue :
\begin{lstlisting}[style=commandstyle]
export LANG=en_US.UTF-8
\end{lstlisting}

\item Vérifier les paramètres des locales :
\begin{lstlisting}[style=commandstyle]
locale
\end{lstlisting}
\end{enumerate}

\section{Ajout des dépôts nécessaires et installation des prérequis}

\begin{enumerate}
\item Installer les outils nécessaires pour gérer les dépôts logiciels :
\begin{lstlisting}[style=commandstyle]
sudo apt install software-properties-common
\end{lstlisting}

\item Ajouter le dépôt "universe" :
\begin{lstlisting}[style=commandstyle]
sudo add-apt-repository universe
\end{lstlisting}

\item Mettre à jour les paquets et installer \texttt{curl} :
\begin{lstlisting}[style=commandstyle]
sudo apt update && sudo apt install curl -y
\end{lstlisting}

\item Télécharger la clé ROS2 :
\begin{lstlisting}[style=commandstyle]
sudo curl -sSL https://raw.githubusercontent.com/ros/rosdistro/master/ros.key -o /usr/share/keyrings/ros-archive-keyring.gpg
\end{lstlisting}

\noindent \underline{Note :} la flèche dans la commande ci-dessous permet de montrer le retour à la ligne involontaire, du à l'écriture de ce tutoriel sur un pdf. Elle est à supprimer dans la commande finale. Cette note s'applique à toutes les commandes du présent document.

\item Ajouter le dépôt ROS2 à la liste des sources :
\begin{lstlisting}[style=commandstyle]
echo "deb [arch=$(dpkg --print-architecture) signed-by=/usr/share/keyrings/ros-archive-keyring.gpg] http://packages.ros.org/ros2/ubuntu $(. /etc/os-release && echo \jammy) main" | sudo tee /etc/apt/sources.list.d/ros2.list > /dev/null
\end{lstlisting}
\end{enumerate}

\section{Installation de ROS2}

\begin{enumerate}
\item Mettre à jour et améliorer le système :
\begin{lstlisting}[style=commandstyle]
sudo apt update
sudo apt upgrade
\end{lstlisting}

\item Installer les paquets de base de ROS2. Deux options sont disponibles selon l'utilisation prévue de ROS2 :
\begin{itemize}
\item \textbf{Installation Desktop (Recommandée)} : Inclut ROS, RViz, des démonstrations et des tutoriels.
\end{itemize}

\begin{lstlisting}[style=commandstyle]
sudo apt install ros-humble-desktop
\end{lstlisting}
\begin{itemize}

\item \textbf{Installation ROS-Base (Version minimale)} : Inclut les bibliothèques de communication, les paquets de messages et les outils en ligne de commande.
\end{itemize}

\begin{lstlisting}[style=commandstyle]
sudo apt install ros-humble-ros-base
\end{lstlisting}

\item Après avoir choisi entre ROS-Base et ROS-Desktop, ajouter l'environnement ROS2 au chemin système :
\begin{lstlisting}[style=commandstyle]
echo source /opt/ros/humble/setup.bash >> ~/.bashrc
source ~/.bashrc
\end{lstlisting}

\noindent Ensuite, installer colcon :
\begin{lstlisting}[style=commandstyle]
sudo apt install python3-colcon-common-extensions
\end{lstlisting}

\end{enumerate}
